\documentclass[10pt]{article}
\usepackage[margin=0.7in]{geometry}
\usepackage{graphicx,fancyhdr,lastpage,longtable,booktabs,fancyvrb,hyperref}
\hypersetup{hidelinks}
\pagestyle{fancy}\fancyhf{}\fancyfoot[C]{\ifnum\value{page}<10 0\fi\thepage/\pageref{LastPage}}
\setlength{\headheight}{14pt}
\begin{document}
\begin{titlepage}
\thispagestyle{fancy}
\centering\vspace*{1in}
{\Huge SIMPLE-16 CPU Design\par}\vspace{1cm}
{\Large CMPE 220 --- November 28, 2026\par}\vspace{1cm}
{\Large DRAFT --- team names required\par}
\vspace{1cm}
Member 1: [name required]\\Member 2: [name required]\\Member 3: [name required]
\vfill
Repository: \url{https://github.com/manishav26/cpu-design}\\
Branch: \texttt{feat/programs-execution-tracing}\\
Draft includes measured execution evidence; outstanding integration issues appear at the end.
\end{titlepage}
\setcounter{page}{2}
\section*{Architecture and instruction encoding}
\begin{figure}[h]\centering
\includegraphics[width=.94\linewidth,height=.46\textheight,keepaspectratio]{../../docs/cpu_diagram.png}
\caption{Team CPU architecture diagram (Member 1).}
\end{figure}
\begin{table}[h]\centering
\begin{tabular}{ll}\toprule
Property & SIMPLE-16 \\\midrule
Word / instruction & 16 bits / fixed 16 bits\\
Registers & R0--R3, unsigned 16-bit values\\
Text / data & 0x00--0x3F / 0x40--0x7F\\
Heap / stack & 0x80--0xBF / 0xC0--0xFF\\
Stack & Downward; empty SP=0x100; first push at 0xFF\\
Flags & Z: zero; N: bit 15; arithmetic and CMP update both\\
Opcode groups / mnemonics & 16 / 18 (PUSH, POP, RET share F)\\\bottomrule
\end{tabular}\caption{Source-verified architecture summary.}
\end{table}
\clearpage
\begin{figure}[h]\centering
\includegraphics[width=\linewidth,height=.82\textheight,keepaspectratio]{../../docs/encoding_table.png}
\caption{Team instruction encoding table (Member 1).}
\end{figure}
\clearpage
\section*{Program listings and manual translation}
R0 holds each result; RAM[0x40] is input and RAM[0x41] is result.
The laptop executes the C reference; SIMPLE-16 executes only assembly.
\subsection*{Native C reference}
\VerbatimInput[fontsize=\small]{../../programs/factorial.c}
\begin{table}[h]\centering
\begin{tabular}{ll}\toprule
C operation & Manual assembly translation\\\midrule
Function argument n & LOAD R1, [input]\\
acc = 1 & LDI R0, 1\\
while (n != 0) & LDI R2, 0; CMP R1, R2; JZ done\\
acc = acc * n & MUL R0, R1\\
n-- & ADDI R1, -1\\
Repeat condition & JMP loop\\
return acc & STORE R0, [result]; HALT\\\bottomrule
\end{tabular}\caption{C-to-assembly mapping; standalone translation inlines the function into the entry program.}
\end{table}
\clearpage
\input{listings.tex}
\section*{Measured results}
\input{results.tex}
Unsigned 8! is 40320 (0x9D80). Its high bit is set, so N=1 is valid.
The loop tests the counter against zero; it does not use N to compare magnitudes.
\subsection*{Validation output}
\VerbatimInput[fontsize=\small]{tests.txt}
\clearpage
\section*{Complete execution evidence}
Every snapshot contains the same contiguous RAM block, 0x0000--0x0041,
including unused text words. Addresses and values are hexadecimal words.
The executed address was captured before stepping; PC is the resulting PC.
IR contains the fetched instruction. MAR/MDR show the last memory access,
which can be the instruction fetch. No instruction has executed in the initial snapshot.
\input{snapshots.tex}
\clearpage
\section*{Build, run, and contribution}
Prerequisites: C11 compiler, Make; Python 3 for test/report generation.
PDF generation additionally uses Tectonic (LaTeX packages may require an initial download).
\begin{Verbatim}[fontsize=\small]
make
make run-example
./build/factorial
./build/simple16 programs/factorial.asm
make traces
make test
make sanitize
make report
\end{Verbatim}
\begin{table}[h]\centering
\begin{tabular}{lp{.72\linewidth}}\toprule
Member & Contribution allocation (names pending)\\\midrule
1 & CPU architecture, ISA, assembler, instruction encoding and architecture figures.\\
2 & CPU execution, registers, ALU, flags, branches, stack, loader and original CLI.\\
3 & Assembly examples, native factorial reference and manual translation; read-only trace formatting and CLI integration; correctness tests; Makefile; run instructions; generated listings, results, snapshots and report integration.\\\bottomrule
\end{tabular}\caption{Specific team contribution allocation, pending confirmation of names and final review.}
\end{table}
\subsection*{Outstanding integration review}
\begin{itemize}
\item Member 2: widen MUL operands before multiplication; signed integer promotion can overflow for large operands. Widen indirect address sums before checking bounds; current narrowing can wrap an invalid address into RAM.
\item Member 1: assembler accepts branch targets up to 0xFF, beyond executable text; .word 65536 silently truncates; source files lack final newlines.
\item UBSan passes the contributed tests. Combined ASan/UBSan stalled on this host and was stopped; no ASan pass is claimed. Example tests cover the submitted arithmetic programs; they do not certify the CPU across all operand values. See report/KNOWN\_ISSUES.md for reproducible issue probes.
\item Add actual team names and confirm allocations. Team review and integration are still required; this draft is not submission-ready.
\end{itemize}
The trace API accepts multiple inclusive memory ranges so later text/data/heap/stack views can reuse the formatter. Recursive execution for the December assignment is outside this contribution.
\end{document}
